\section{Definition of Terms}
\subsection{Turf House}

A turf house is defined as a structure or building whose roof and walls
are mainly made from turf. Sometimes, these can be mixed with gravel or
stones for extra support. Timber beams and rafters are used as the main
structure of the building and are also used to support and hold the
roof. Turf is also used as a main component to cover the roof, and
slates are used as an underlay. Throughout the years, the material used
for turf houses has not changed drastically, but the form of these
houses has changed and adapted through the ages \cite{ref14}\cite{ref15}.

\subsection{Soil Moisture}

Soil moisture is defined as the amount or quantity of water in the soil
\cite{ref16}\cite{ref17}. It can be further defined as the amount of water held
in between the pores of the soil. These volumes of water may be in
liquid or vaporized form \cite{ref18}. When measuring the amount of water in
a known amount of soil, soil moisture can be expressed as ``\% water by
weight or volume of soil'', or ``inches of water per foot of soil''
\cite{ref19}.

\subsection{Relative Humidity in Soil}

Humidity is defined as the amount of gaseous water or water vapor in the
air. High humidity means there is a lot of moisture in the air and low
humidity indicates dry conditions. Humidity is often described as
absolute humidity (AH) or relative humidity (RH) \cite{ref20}\cite{ref21}.
Absolute humidity is measured in units of g/m\textsuperscript{3} where
it is the measure of water vapor content in air that is not dependent on
temperature \cite{ref22}. Relative humidity is the ratio of the current
amount of water vapor in the air to the maximum amount the air can hold
at a given temperature, expressed as a percentage (\%) \cite{ref23}.

This study will collect data and conduct tests below the soil surface.
In this context, humidity refers to the relative humidity of the air
within the soil\textquotesingle s pores and gaps, which will be used as
data for the study.

\subsection{Wireless Sensor Network (WSN)}

A wireless sensor network (WSN) is a collection of spatially distributed
sensor nodes that measure and communicate data wirelessly using protocols
such as radio, Wi-Fi, Bluetooth, and LoRa. A sensor node may contain one
or more sensors, depending on the application, and can communicate with a
central system or with neighboring nodes in a mesh network
\cite{jung2012clusterhead,xu2014borderrouter}. WSNs are used in
environmental monitoring, industrial equipment monitoring, healthcare,
agriculture, and other distributed sensing applications. Their low-power
design makes them suitable for remote and harsh environments.

\subsection{Internet of Things (IoT)}

The Internet of Things (IoT) refers to interconnected physical devices
that collect, exchange, and process data through communication networks
and cloud services \cite{aws2024iot}. Examples include wearable devices,
tracking systems, environmental sensors, and smart infrastructure. In this
context, a WSN provides the sensing and communication layer, while the IoT
system can transmit the collected measurements to a network or cloud
platform for storage, visualization, and analysis
\cite{karthik2021iot,jebri2015iot}.

\subsection{Long Range (LoRa)}

Long Range (LoRa) is a radio-frequency modulation technology designed for
low-power wide-area network (LPWAN) applications. It supports long-range
communication with relatively low energy consumption, making it suitable
for many IoT applications and outdoor monitoring systems
\cite{dambal2019lora}. Communication range depends on factors such as
line of sight, antenna design, frequency, spreading factor, bandwidth,
transmission power, and local regulations. The sensor nodes in this study
use a 433 MHz SX1278-based LoRa module in the applicable industrial,
scientific, and medical (ISM) band. The use of 433 MHz, 868 MHz, or
915 MHz bands depends on regional regulations and deployment requirements
\cite{camarillo2022lorawan,yahya2023antenna}.

\subsection{In-Situ Sensors}

In-situ sensors measure and monitor conditions directly at the location of
interest. They provide measurements without requiring repeated manual
sampling or removal of material from the monitored environment. In-situ
sensors are useful in environmental monitoring because they can collect
data continuously from remote locations. They may include temperature,
air-quality, soil-moisture, and groundwater sensors and can be integrated
with processing and wireless communication hardware
\cite{barnhart2010groundwater}.

\subsection{Low-Power Systems}

Low-power systems are designed to operate with minimal energy consumption
over extended periods, often using batteries or harvested energy. These
systems commonly operate intermittently and use hardware and software
optimization to reduce their duty cycle. Sleep modes disable non-essential
functions while preserving the control functions needed to wake the device
at scheduled intervals. This approach is important for battery-powered IoT
devices deployed in remote areas \cite{skuta2020power, montoya2018sward}.

\subsection{Cloud-Based Data Platforms}

Cloud-based data platforms provide online infrastructure for storing,
managing, processing, and analyzing data. They allow users to access data
remotely without maintaining local high-capacity storage or computing
hardware. Cloud services are commonly used in IoT systems for device
management, data processing, visualization, and analytics. Examples include
Amazon Web Services, Google Cloud, and Microsoft Azure. In this thesis,
ThingSpeak provides the cloud-based storage and visualization service for
the collected sensor measurements \cite{sun2011cloud}.

\section{Background of the Study}
Icelandic turf houses are important cultural and architectural heritage structures whose turf construction requires continued maintenance. The literature supplied for this thesis describes the long history of turf construction in Iceland and the need to understand turf layers and their environmental conditions. Moisture is particularly relevant because turf and timber components are susceptible to deterioration, while the conservation of heritage structures requires maintenance approaches that do not unnecessarily disturb the site.

The current system design addresses this monitoring problem by placing low-power sensor nodes beneath turf roofs. The methodology specifies that nodes collect soil moisture, soil temperature, air temperature, and relative humidity, while transmitting measurements through a LoRa mesh network to a base station and cloud database. The nodes are designed to remain visually unobtrusive and to operate for extended periods with limited maintenance.

\section{Statement of the Problem}
Monitoring environmental conditions beneath turf roofs presents several engineering constraints. Sensor nodes must be sufficiently small and unobtrusive to avoid affecting the appearance of a heritage site, while also operating in an environment characterized by high moisture, freezing and thawing cycles, and possible snow coverage. The system must provide reliable wireless communication, conserve battery energy, retain measurements when transmissions are missed, and make the collected data available remotely.

The project therefore addresses the need for a practical monitoring system that combines low-cost sensing hardware, low-power operation, long-range wireless communication, local data backup, and cloud-based access in a single deployable platform.

\section{Objectives}
The study aims to design, implement, and evaluate a low-power LoRa wireless sensor mesh network for monitoring soil and environmental conditions beneath Icelandic turf roofs. The specific objectives are to:
\begin{enumerate}
\item measure soil moisture, soil temperature, air temperature, relative humidity, and battery voltage using distributed sensor nodes;
\item develop a LoRa-based mesh communication system capable of reliably transferring measurements between sensor nodes and a base station;
\item implement low-power operating modes suitable for long-duration battery-powered deployment;
\item provide local data backup so measurements are retained when wireless transmissions are missed;
\item calibrate the low-cost capacitive soil moisture sensor against reference measurements; and
\item transmit collected measurements to a cloud platform for remote access, visualization, and analysis.
\end{enumerate}

\section{Significance of the Study}
The project provides a low-cost engineering approach to continuous monitoring of environmental conditions on Icelandic turf roofs. The system is intended to support researchers and conservation personnel by providing measurements without requiring frequent manual sampling or intrusive installation. The use of off-the-shelf components and a modular mesh architecture also provides a basis for scaling the monitoring approach to additional turf-roof sites.

\section{Scope and Limitations}
The work focuses on a low-power LoRa-based wireless sensor mesh network using distributed sensor nodes and a cloud-connected base station. The implemented nodes measure soil moisture, soil temperature, air temperature, relative humidity, and battery voltage. The methodology specifies burial of sensor nodes at approximately 11 cm, weatherproof enclosures, local SD-card data backup, and a non-destructive installation approach.

The evaluation includes laboratory power-consumption testing, preliminary and intermediate field deployments around the University of Iceland, LoRa range testing around Vellir in Hafnarfjörður, long-term deployment at the Keldur turf-house site, soil moisture calibration, and ThingSpeak cloud-upload testing. The documented results also show limitations, including antenna-related packet loss during early testing, sensor-calibration issues during intermediate testing, and hardware failures affecting two nodes near the end of the long-term deployment. The calibration results cover the 0--40\% moisture range and therefore do not establish performance beyond that range.